\documentclass{beamer} 
%\documentclass[handout]{beamer} 

% Michael Maier, 2016.
% CC-0

\usepackage[utf8]{inputenc}
\usepackage[ngerman]{babel}

\title{Android Googlefrei} 
\author{Michael Maier \textless Michael.Maier@mailbox.org\textgreater} 
\date{26. April 2016} 

\usetheme{Antibes}

\hypersetup{colorlinks=true,urlcolor=blue,linkcolor=white}

%\usebackgroundtemplatei{
%\includegraphics[width=\paperwidth,
%height=0.8\paperheight]{mag_map.png}
%}

\begin{document}

%\maketitle

\begin{frame} 


\begin{figure}
  \centering
  \includegraphics[width=.8\textwidth]{titel.pdf}
\end{figure}

\begin{center}
\Huge{Android Googlefrei\\}
\end{center}

%\begin{center}
%\Large{\emph{Selbsterfasste Daten als Grundlage für Navigation}}
%\end{center}

\end{frame}


\section{Einleitung}


\begin{frame}{Vorstellung}

  \begin{itemize}
    \item Michael Maier \textless \href{mailto:Michael.Maier@mailbox.org}{Michael.Maier@mailbox.org}\textgreater
    \item Student an der TU Graz (Telematik)
\vspace{0.3cm}
    \item Linux-User (Debian/grml) seit 2004
    \item Organisiere Grazer Linuxtage seit 2011 mit
    \item OpenStreetMap als Hobby seit Juli 2010
    \item Leite den Grazer OSM-Stammtisch seit Mai 2011
\vspace{0.3cm}
    \item Freiberuflich OSM-Aufträge und Consulting
    \item "`Regular"' bei der Cryptoparty Graz
  \end{itemize}
\end{frame}

\subsection{Werdegang}
\begin{frame}{Meine Telefon-Geschichte}

\begin{columns}[c]
  \begin{column}[T]{.72\textwidth}
    \begin{itemize}
      \item War lange Zeit Smartphone-Verweigerer\dots
      \item 10 Jahre Nokia 6310i
        \begin{itemize}
          \item SMS, Kalender-, Kontakte-Sync via Bluetooth
          \item Trauere um Profile...
        \end{itemize}
      \item Unfreie Systeme kommen nicht mehr in Frage.
    \end{itemize}
  \end{column}
  \begin{column}[T]{.28\textwidth}
  \vspace{-2.5cm}
    \includegraphics[width=2cm]{nokia.png} \\
    {\TINY \vspace{.05cm}CC-BY-SA 3.0 \\ \vspace{-.25cm} WP:User Martinz6492}
 \end{column}
\end{columns}

\pause
\begin{columns}[c]
  \begin{column}[T]{.88\textwidth}
    \vspace{.5cm}
    \begin{itemize}
      \item "`Smart"' angefangen mit OpenMoko ... Da läuft auch Debian
      \item War ein guter Mobilcomputer, als Telefon leider nicht praktikabel (Sprachqualität)
    \end{itemize}
    \hfill{\TINY CC-BY-SA 3.0 WP:User Nevetsjc}
  \end{column}
  \begin{column}[T]{.11\textwidth}
    \includegraphics[width=2cm]{neo.png} 
   \end{column}
\end{columns}

\end{frame}

  

\begin{frame}{Android lange mit Sorge beobachtet}

Dann kam das Fairphone
\begin{itemize}
  \item Es kommt gerootet und ohne Google
  \item Und es ist auch noch ein wenig Fairer als andere

\pause

  \item Android 4.2.2 "`Jelly Bean"' \hfill \includegraphics[width=2cm]{jellybean.png}
\end{itemize}
  
  Darauf basieren meine Erfahrungen - jedes Telefon ist anders!
\pause

\begin{itemize}
  \item Inzwischen: Fairphone2 
  \begin{itemize}
    \item leider noch nicht ohne Google (rooten möglich)
    \item Android 5.1
  \end{itemize}
\end{itemize}

\end{frame}

\begin{frame}{Voraussetzungen}

  Ausgangssituation:
\begin{itemize}
  \item Gerootetes Handy (4.2.2) ohne Play-Store/Services
  \item OwnCloud oder Vertrauenswürdiger WebDAV-Anbieter
    \begin{itemize}
      \item z.B. mailbox.org
    \end{itemize}

\end{itemize}

  \vspace{1cm}
Ich sag es gleich - ich hab keinen Google-Account, und kann darum keinen exakten Vergleich zwischen den Apps machen
  \vspace{1cm}

\pause

\begin{itemize}
  \item Einige Monate parallel zu 6310i als Organizer verwendet, bis alles so eingerichtet wie vom Nokia gewohnt.
\end{itemize}

Damit ihr keine 3 Monate mehr braucht, dafür dieser Talk!

\end{frame}

\section{Apps-Store}

\begin{frame}{App-Store: F-Droid}

\begin{columns}[c]
  \begin{column}[T]{.7\textwidth}
    \begin{itemize}
      \item  Ohne Play Store... 
      \item Muß man die Installation von "`Fremd"'-Apps erlauben

        \pause

      \item Unsicher? \pause Nicht unsicherer als Play Store!
        \begin{itemize}
          \item Betrüger treiben mittels Android-VMs die Downloadzahlen und Bewertungen in die Höhe
        \end{itemize}
    \end{itemize}
  \end{column}
  \begin{column}[T]{.3\textwidth}
    \includegraphics[width=3.5cm]{fremdapps.png} 
   \end{column}
\end{columns}

\pause

\begin{columns}[c]
  \begin{column}[T]{.7\textwidth}

    \vspace{-2cm}
      \begin{itemize}
      \item 100\,\% Schutz vor Viren gibts keinen, aber:
        \begin{itemize}
          \item F-Droid baut aus den Sourcen
          \item  Ihr könnt euch anschauen welche Entwickler dahinterstecken.
        \end{itemize}
      \end{itemize}
   \end{column}
  \begin{column}[T]{.3\textwidth}
    \vspace{-6.1cm}
    \includegraphics[width=3.5cm]{fdroid.png} 
   \end{column}
\end{columns}

\end{frame}

\section{Datenschutz}

\begin{frame}{Mehr Datenschutz: XPrivacy}

\begin{columns}[c]
  \begin{column}[T]{.7\textwidth}
    \begin{itemize}
      \item Zuerst Xposed-Framework installieren: \href{http://xposed.info}{xposed.info}
      \item Rebooten
      \item XPrivacy übers Framework installieren
    \end{itemize}
  \end{column}
  \begin{column}[T]{.3\textwidth}
    \includegraphics[width=3.5cm]{xposed.png} 
   \end{column}
\end{columns}

        \pause

\begin{columns}[c]
  \begin{column}[T]{.7\textwidth}

    \vspace{-2cm}
      \begin{itemize}
        \item Apps Rechte wieder wegnehmen
        \item eure Privaten Daten verschleiern
      \end{itemize}
   \end{column}
  \begin{column}[T]{.3\textwidth}
    \vspace{-6.1cm}
    \includegraphics[width=3.5cm]{xprivacy.png} 
   \end{column}
\end{columns}

        \pause

\begin{columns}[c]
  \begin{column}[T]{.7\textwidth}
    \vspace{-1cm}
      \begin{itemize}
        \item 100\,\%-Schutz gibt es nicht, an viele Daten kommen die Apps auch auf anderem Wege
      \end{itemize}
   \end{column}
  \begin{column}[T]{.3\textwidth}
   \end{column}
\end{columns}

\end{frame}

\begin{frame}{Datenschutz: Sonstige Einstellungen}

\begin{columns}[c]
  \begin{column}[T]{.7\textwidth}
    WLAN $\Longrightarrow$ Erweitert $\Longrightarrow$ Netzwerkhinweis deaktivieren
    \begin{itemize}
      \item Sonst versucht sich Android mit allen bekannten Netzen zu verbinden und Broadcasted die bisher besuchten SSIDs
    \end{itemize}
  \end{column}
  \begin{column}[T]{.3\textwidth}
    \includegraphics[width=3.5cm]{wlan.png} 
   \end{column}
\end{columns}

        \pause

\begin{columns}[c]
  \begin{column}[T]{.7\textwidth}

    \vspace{-2cm}
    Standordzugriff (GPS) $\Longrightarrow$ Standort per WLAN und Mobilfunknetz deaktivieren
      \begin{itemize}
        \item Telefoniert nach Hause und meldet Google, in welchen WLAN / Mobilfunk-Zelle ich bin
      \end{itemize}
   \end{column}
  \begin{column}[T]{.3\textwidth}
    \vspace{-6.1cm}
    \includegraphics[width=3.5cm]{standort.png} 
   \end{column}
\end{columns}


\end{frame}

\section{Weg von der "`Other People's Computers"' Cloud}

\subsection{Damit private Daten privat bleiben}

\begin{frame}{Kontakte, Termine}
  

\begin{columns}[c]
  \begin{column}[T]{.7\textwidth}
    App "`DavDroid"' \includegraphics[width=0.5cm]{dav.png}

    \begin{itemize}
      \item Verbindung zur Owncloud
      \item Legt ein neues Benutzerkonto am Gerät an
    \vspace{0.5cm}
      \item CardDav: Kontakte
    \end{itemize}
  \end{column}
  \begin{column}[T]{.3\textwidth}
    \includegraphics[width=3.5cm]{carddav.png} 
   \end{column}
\end{columns}

        \pause

\begin{columns}[c]
  \begin{column}[T]{.7\textwidth}

    \vspace{-2cm}
      \begin{itemize}
        \item CalDav: Kalender
        \begin{itemize}
          \item Mit extra App "`PoliteDroid"': Ton aus bei "`busy"' Terminen \includegraphics[width=0.5cm]{polite.png}
        \end{itemize}
      \item Zusätzlich: App "`ICSdroid"' \includegraphics[width=0.5cm]{ics.png}
      \end{itemize}
   \end{column}
  \begin{column}[T]{.3\textwidth}
    \vspace{-6.1cm}
    \includegraphics[width=3.5cm]{caldav.png} 
   \end{column}
\end{columns}

\end{frame}
  
   

\begin{frame}{Cloud-Speicher}

\begin{columns}[c]
  \begin{column}[T]{.7\textwidth}

    \vspace{0.5cm}
    Dateisynchronisierung: ownCloud
    \vspace{0.5cm}

    \begin{itemize}
      \item Der Filesync-Client zum Vortrag um 10:00.
    \end{itemize}
  \end{column}
  \begin{column}[T]{.3\textwidth}
    \includegraphics[width=3.5cm]{owncloud.png} 
   \end{column}
\end{columns}

        \pause

\begin{columns}[c]
  \begin{column}[T]{.7\textwidth}

    \vspace{-2cm}
      \begin{itemize}
        \item Fotos sofort in die Cloud schieben: check.
      \end{itemize}
   \end{column}
  \begin{column}[T]{.3\textwidth}
    \vspace{-6.1cm}
    \includegraphics[width=3.5cm]{owncloud-upload.png} 
   \end{column}
\end{columns}

\end{frame}


\begin{frame}{SMS-Backup}
\begin{columns}[c]
  \begin{column}[T]{.7\textwidth}
      App "`SMS Backup+"'
    \begin{itemize}
      \item Sichert SMS (empfangen/gesendet) auf einen IMAP-Server
      \item Ebenso die Anruflisten-Logdatei
    
    \end{itemize}
  \end{column}
  \begin{column}[T]{.3\textwidth}
    \includegraphics[width=3.5cm]{sms.png} 
   \end{column}
\end{columns}
\end{frame}

\subsection{Google-Dienste ersetzen}

\begin{frame}{Karten \& Navi}

\begin{columns}[c]
  \begin{column}[T]{.7\textwidth}
    OpenStreetMap statt Google Maps!

    \begin{itemize}
      \item OsmAnd (OsmAnd\~{} in F-Droid)  \includegraphics[width=0.5cm]{osma.png}
      \begin{itemize}
        \item Komplett Offline
        \item POI-Suche (zB alle Fleischereien)
        \item Fußgänger/Fahrrad/Autonavi mit Spurassistent
      \end{itemize}
    \end{itemize}
  \end{column}
  \begin{column}[T]{.3\textwidth}
    \includegraphics[width=3.5cm]{osmand.png} 
   \end{column}
\end{columns}

        \pause

\begin{columns}[c]
  \begin{column}[T]{.7\textwidth}

    \vspace{-2.4cm}
      \begin{itemize}
        \item Reicht nicht? Es gibts \textgreater 100 weitere OSM-Apps für Android, siehe OSM-Wiki: \url{http://wiki.osm.org/Android}
      \end{itemize}
   \end{column}
  \begin{column}[T]{.3\textwidth}
   \end{column}
\end{columns}
\end{frame}


\begin{frame}{VoIP: statt Hangouts/Skype}

\begin{columns}[c]
  \begin{column}[T]{.7\textwidth}
  
    Mumble (Murmur) als VoIP-Server, die App "`Plumble"' als Android-Client \includegraphics[width=0.5cm]{plumblei.png}

    \begin{itemize}
      \item Audiokonferenzen für 100+ Leute (oder zu 2. \dots)
      \item Sichere End2End-Verschlüsselung
      \item Server- und Clientzertifikate (optional)
    \end{itemize}

    
  \end{column}
  \begin{column}[T]{.3\textwidth}
    \includegraphics[width=3.5cm]{plumble.png} 
   \end{column}
\end{columns}

\end{frame}


\begin{frame}{Instant-Messenger}

\begin{columns}[c]
  \begin{column}[T]{.7\textwidth}
    \vspace{-0.2cm}
  \includegraphics[width=0.5cm]{sig.png} ~  LibreSignal statt Whats?App
    

    \begin{itemize}
      \item WhatsApp nutzt inzwischen Signal-Verschlüsselung - gleich beim Original bleiben :-) 
        \item LibreSignal wurde ohne Google-Play-Services kompiliert
        \item Keine Abstriche bei Text und Bildern
    \end{itemize}
  \end{column}
  \begin{column}[T]{.3\textwidth}
    \includegraphics[width=3.5cm]{signal.png} 
   \end{column}
\end{columns}
        \pause

\begin{columns}[c]
  \begin{column}[T]{.7\textwidth}

    \vspace{-2.4cm}
      \begin{itemize}
        \item Telefonieren hat bei mir mit LibreSignal nicht funktioniert.
        \item Es wird die Kontaktliste ungefragt zum Server gesendet, wenn Zugriff gestattet ($\Longrightarrow$ blocken per XPrivacy)
        \begin{itemize}
          \item Ansonsten nur Telefonnumern statt Namen
        \end{itemize}
      \end{itemize}
   \end{column}
  \begin{column}[T]{.3\textwidth}
   \end{column}
\end{columns}

\end{frame}

\subsection{Nützliches aus dem F-Droid Store}

\begin{frame}{Ich kann nicht mehr leben ohne:}

\begin{itemize}
  \item \includegraphics[width=0.5cm]{bar.png} "`Barcode Scanner"' (ohne Werbung!) 
  \item \includegraphics[width=0.5cm]{clip.png} "`ClipStack"' - mehr als eine Zwischenablage  
  \item \includegraphics[width=0.5cm]{keyb.png} "`Hacker's Keyboard"' (C-A,C-X anyone?) 
  \item \includegraphics[width=0.5cm]{fennec.png} "`Fennec"' (Firefox)
  \item \includegraphics[width=0.5cm]{tux.png} "`GLT-Begleiter"'
  \item \includegraphics[width=0.5cm]{ovpn.png} "`OpenVPN"'
  \item \includegraphics[width=0.5cm]{red.png} "`Red Moon"' - f.lux fürs Handy
  \item \includegraphics[width=0.5cm]{night.png} "`Silent Night"' ((root)) - Nacht: Flugmodus
  \item \includegraphics[width=0.5cm]{wifi.png} "`Smarter WiFi-Manager"' - WLAN abh. von Funkzelle
\end{itemize}

\end{frame}


\section{Ende}

\begin{frame}{Hilfe?}

  Zu kompliziert? Zu schnell?
    \vspace{1cm}

  \pause

  Kommt zur Cryptoparty, hier werden sie geholfen!

  \begin{itemize}
    \item Jeden 1. Montag im Monat im Spekral Graz (Lendkai 45)
    \item z.B. übermorgen Montag, 2.5.2016
  \end{itemize}

    \vspace{1cm}
      Nicht vergessen: Vorträge bewerten :-)

\end{frame}

\begin{frame}{Vielen Dank für die Aufmerksamkeit!}

  Folien für die Grazer Linuxtage, 30.4.2016
\vspace{1cm}

Erstellt mittels \LaTeX Beamer, Quelltext: \href{https://github.com/species/android-googlefrei}{Github/species/android-googlefrei}.
\vspace{1cm}

\href{mailto:michael.maier@mailbox.org}{Michael Maier}

Twitter: \href{https://twitter.com/osmgraz}{@osmgraz}
\vspace{1cm}

Folien unter: \includegraphics[width=1cm]{cc-by-sa.pdf}. 


\end{frame}



\end{document}
